\begin{proof}
    Let $U \in \CatUnitary( d, d )$ and define $F( - ) = U \odot ( - )$.
    Let $\{ \alpha_j \}_{j=1}^d$ denote the eigenvalues of $\Log( U ) / (i\pi)$ with orthonormal eigenbasis $\{ b_j \}_{j=1}^d$.
    There are four properties to verify.
    \begin{enumerate}
    \item $F( 1_n ) = \exp( \Log( U ) \otimes \Log( 1_n ) / ( i\pi ) ) = \exp( \Log( U ) \otimes \mathbb{O}_n / ( i\pi ) ) = \exp( \mathbb{O}_{dn} ) = \mathbb{I}_{dn}$ for $n \in \mathbb{N}$.
    \item Let $n \in \mathbb{N}$ and $V \in \Mor( \CatUnitary )$.
          Then the following equation holds.
          \begin{align*}
              U \odot ( V \otimes \mathbb{I}_n )
              &=
              \sum_{j=1}^d \ket{b_j} \bra{b_j} \otimes ( V \otimes \mathbb{I}_n )^{\alpha_j}
              &
              ( \text{by Cor.~\ref{Thm:GenCtrlNF}} )
              \\
              &=
              \sum_{j=1}^d \ket{b_j} \bra{b_j} \otimes \exp( \alpha_j \Log( V \otimes \mathbb{I}_n ) )
              \\
              &=
              \sum_{j=1}^d \ket{b_j} \bra{b_j} \otimes \exp( \alpha_j \Log( V ) \otimes \mathbb{I}_n )
              \\
              &=
              \sum_{j=1}^d \ket{b_j} \bra{b_j} \otimes \exp( \alpha_j \Log( V ) ) \otimes \mathbb{I}_n
              \\
              &=
              \left( \sum_{j=1}^d \ket{b_j} \bra{b_j} \otimes \exp( \alpha_j \Log( V ) ) \right) \otimes \mathbb{I}_n
              \\
              &=
              \left( \sum_{j=1}^d \ket{b_j} \bra{b_j} \otimes V^{\alpha_j} \right) \otimes \mathbb{I}_n
              =
              F( V ) \otimes \mathbb{I}_n
              &
              ( \text{by Cor.~\ref{Thm:GenCtrlNF}} )
          \end{align*}
    \item Let $V \in \CatUnitary( n, n )$.
          Since $\sigma_{d,d}$ is a monoidal symmetry, then $\sigma_{d,d}$ is self-adjoint with $\sigma_{d,d} \cdot ( \Log( U ) \otimes \Log( U ) ) \cdot \sigma_{d,d} = \Log( U ) \otimes \Log( U )$.
          Then the following equation holds where $H = \Log( V )$.
          \begin{align*}
              F( F( V ) )
              &=
              U \odot ( U \odot V )
              \\
              &=
              \exp( \Log( U ) \otimes \Log( U ) \otimes H / (i\pi)^2 )
              &
              ( \text{by Thm.~\ref{Thm:ControlAssoc}} )
              \\
              &=
              \exp( ( \sigma_{d,d} \cdot ( \Log( U ) \otimes \Log( U ) ) \cdot \sigma_{d,d} ) \otimes H / (i\pi)^2 )
              \\
              &=
              \exp( ( \sigma_{d,d} \otimes \mathbb{I}_n ) \cdot ( \Log( U ) \otimes \Log( U ) \otimes H / (i\pi)^2 ) \cdot ( \sigma_{d,d} \otimes \mathbb{I}_n ) )
              \\
              &=
              ( \sigma_{d,d} \otimes \mathbb{I}_n ) \cdot \exp( \Log( U ) \otimes \Log( U ) \otimes H / (i\pi)^2 ) \cdot ( \sigma_{d,d} \otimes \mathbb{I}_n )
              &
              ( \text{by~\cref{Eqn:UnitaryConj}} )
              \\
              &=
              ( \sigma_{d,d} \otimes \mathbb{I}_n ) \cdot ( U \odot ( U \odot V ) ) \cdot ( \sigma_{d,d} \otimes \mathbb{I}_n )
              &
              ( \text{by Thm.~\ref{Thm:ControlAssoc}} )
              \\
              &=
              ( \sigma_{d,d} \otimes \mathbb{I}_n ) \cdot F( F( V ) ) \cdot ( \sigma_{d,d} \otimes \mathbb{I}_n )
          \end{align*}
          Then $F( F( V ) ) = ( \sigma_{d,d} \otimes \mathbb{I}_n ) \cdot F( F( V ) ) \cdot ( \sigma_{d,d} \otimes \mathbb{I}_n )$.
          Since $\sigma_{d,d} \otimes \mathbb{I}_n$ is also self-adjoint, then $( \sigma_{d,d} \otimes \mathbb{I}_n ) \cdot F( F( V ) ) = F( F( V ) ) \cdot ( \sigma_{d,d} \otimes \mathbb{I}_n )$.
    \item If $V \in \CatUnitary( n, n )$ and $P \in \CatUnitary( n, n )$, then the following equation holds.
          \begin{align*}
              F( P^\dagger \circ V \circ P )
              &=
              \exp( \Log( U ) \otimes \Log( P^\dagger \circ V \circ P ) / (i\pi) )
              \\
              &=
              \exp( \Log( U ) \otimes ( P^\dagger \circ \Log( V ) \circ P ) / (i\pi) )
              &
              ( \text{by~\cref{Eqn:UnitaryConj}} )
              \\
              &=
              \exp( ( \mathbb{I}_d \otimes P^\dagger ) \circ ( \Log( U ) \otimes \Log( V ) / (i\pi) ) \circ ( \mathbb{I}_d \otimes P ) )
              \\
              &=
              ( \mathbb{I}_d \otimes P^\dagger ) \circ \exp( \Log( U ) \otimes \Log( V ) / (i\pi) ) \circ ( \mathbb{I}_d \otimes P )
              &
              ( \text{by~\cref{Eqn:UnitaryConj}} )
              \\
              &=
              ( \mathbb{I}_d \otimes P^\dagger ) \circ F( V ) \circ ( \mathbb{I}_d \otimes P )
          \end{align*}
    \end{enumerate}

    \noindent
    \textbf{Functoriality}.
    %
    Next, it must be shown that $F( - )$ is an ordinary functor if and only if $U$ is Hermitian.
    As a preliminary result, a general equation will be computed for $F( V \circ W)$.
    Let $V, W \in \CatUnitary( n, n )$.
    Then by \cref{Thm:GenCtrlNF}, the following equations hold.
    \begin{align*}
        F( V ) &= \sum_{j=1}^{d} \ket{b_j} \bra{b_j} \otimes V^{\alpha_j}
        &
        F( W ) &= \sum_{j=1}^{d} \ket{b_j} \bra{b_j} \otimes W^{\alpha_j}
    \end{align*}
    Since $\{ b_j \}_{j=1}^d$ is orthonormal, then $\braket{b_j|b_k} = \delta_{j,k}$ where $\delta$ is the Kronecker delta function.
    Then the following equation holds by the bilinearity of $( \otimes )$.
    \begin{equation*}
        F( V ) \circ F( W )
        =
        \sum_{j=1}^{d} \sum_{k=1}^{d} \delta_{j,k} \ket{b_j} \bra{b_k} \otimes ( V^{\alpha_j} \circ W^{\alpha_k} )
        =
        \sum_{j=1}^{d} \ket{b_j} \bra{b_j} \otimes ( V^{\alpha_j} \circ W^{\alpha_j} )
    \end{equation*}
    Since $V$ and $W$ were arbitrary, then this equation holds in general.
    Now it can be shown that $F( - )$ is a functor if and only if $U$ is Hermitian.
    \begin{itemize}
    \item Assume that $F( - )$ is a functor.
          Let $j \in \{ 1, 2, \ldots, d \}$.
          Observe that the following equations hold where $M = \ket{b_j} \otimes \mathbb{I}_2$.
          \begin{align*}
              M^\dagger \circ F( -Z ) \circ M
              &=
              \sum_{k=1}^{d} \delta_{j,k} \delta_{k,j} ( -Z )^{\alpha_k}
              =
              ( -Z )^{\alpha_j}
              \\
              M^\dagger \circ F( -\mathbb{I}_2 ) \circ F( Z ) \circ M
              &=
              \sum_{k=1}^{d} \delta_{j,k} \delta_{k,j} ( -\mathbb{I}_2 )^{\alpha_k} Z^{\alpha_k}
              =
              ( -\mathbb{I}_2 )^{\alpha_j} \circ Z^{\alpha_j}
          \end{align*}
          Since $F( - )$ is a functor, then the following equation holds.
          \begin{equation*}
              ( -Z )^{\alpha_j}
              =
              M \circ F( -Z ) \circ M
              =
              M \circ F( -\mathbb{I}_2 ) \circ F( Z ) \circ M
              =
              ( -\mathbb{I}_2 )^{\alpha_j} \circ Z^{\alpha_j}
              =
              e^{i\alpha_j\pi} Z^{\alpha_j}
          \end{equation*}
          Since $\Log( -Z ) = i \ket{0}\bra{0} \pi$ and $\Log( Z ) = i \ket{1}\bra{1} \pi$, then the following equation holds.
          \begin{align*}
            ( -Z )^{\alpha_j}
            &=
            \exp( i\alpha_j\pi \ket{0} \bra{0} )
            =
            \begin{bsmallmatrix}
                e^{i\alpha_j\pi} & 0 \\
                0 & 1
            \end{bsmallmatrix}
            &
            Z^{\alpha_j}
            &=
            \exp( i\alpha_j\pi \ket{1} \bra{1} )
            =
            \begin{bsmallmatrix}
                1 & 0 \\
                0 & e^{i\alpha_j\pi}
            \end{bsmallmatrix}
          \end{align*}
          Then the following equation also holds.
          \begin{equation*}
              1
              =
              \braket{1|1}
              =
              \bra{1} ( -Z )^{\alpha_j} \ket{1}
              =
              e^{i\alpha_j\pi} \bra{1} Z^{\alpha_j} \ket{1}
              =
              e^{i\alpha_j\pi} \left( e^{i\alpha_j\pi} \braket{1|1} \right)
              =
              e^{i\alpha_j(2\pi)}
          \end{equation*}
          Then $\alpha_j(2\pi)$ is a multiple of $2\pi$.
          Then $\alpha_j$ is an integer.
          Then $e^{i\alpha_j\pi}$ is either $-1$ or $1$.
          Since $j$ was arbitrary, then all eigenvalues of $U$ are either $-1$ or $1$.
          Since $U$ is unitary, then this means $U$ is also Hermitian.
    \item Assume that $U$ is Hermitian.
          Since $U = \exp( iH\pi )$, then the eigenvalues of $U$ are $\{ e^{i\alpha_j\pi} \}_{j=1}^{d}$.
          Let $j \in \{ 1, 2, \ldots, d \}$.
          Since $U$ is Hermitian and unitary, then $e^{i\alpha_j\pi}$ is either $1$ or $-1$.
          Then $\alpha_j\pi$ is a multiple of $\pi$.
          Since $\alpha_j \in (-1, 1]$, then $\alpha_j \in \{ 0, 1 \}$.
          Then either $\alpha_j = 0$ and $V^{\alpha_j} \circ W^{\alpha_j} = \mathbb{I}_m = ( V \circ W )^{\alpha_j}$ or $\alpha_j = 1$ and $V^{\alpha_j} \circ W^{\alpha_j} = V \circ W = ( V \circ W )^{\alpha_j}$.
          Since $j$ was arbitrary, then $V^{\alpha_j} \circ W^{\alpha_j} = ( V \circ W )^{\alpha_j}$.
          Then the following equation holds.
          \begin{equation*}
              F( V ) \circ F( W )
              =
              \sum_{j=1}^{d} \ket{b_j} \bra{b_j} \otimes ( V^{\alpha_j} \circ W^{\alpha_j} )
              =
              \sum_{j=1}^{d} \ket{b_j} \bra{b_j} \otimes ( V \circ W )^{\alpha_j}
              =
              F( V \circ W )
          \end{equation*}
          Since $V$ and $W$ were arbitrary, then $F( - )$ preserves composition.
          It was already shown in (\textbf{C1}) that $F( - )$ preserves identities, so $F( - )$ is a functor.
    \end{itemize}
    Then $F( - )$ is a functor if and only if $U$ is Hermitian.

    \vspace{1em}
    \noindent
    \textbf{Control Functoriality}.
    %
    By definition, $( \CatUnitary_d )_\mathbf{endo} = \CatUnitary_d$.
    Notice that for each $k \in \mathbb{N}$, $\Ob( F )( d^k ) = d d^k = d^{k+1}$.
    Then there exists a functor $C: \CatUnitary_d \to \CatUnitary_d$ such that $C( V: k \to k ) = F( V ) \in \CatUnitary( d^{k+1}, d^{k+1} ) = \CatUnitary_d( k + 1, k + 1 )$.
    It can then be shown that $C( - )$ is a conjugated control functor.
    \begin{itemize}
    \item If $V \in \CatUnitary_d( k, k )$, then $C( V \otimes 1 ) = F( V \otimes \mathbb{I}_d ) = F( V ) \otimes \mathbb{I}_d = C( V ) \otimes 1$ by (\textbf{C2}).
    \item Let $V \in \CatUnitary_d( k, k )$ and write $n = d^k$.
          It follows from (\textbf{C3}) that $( \sigma \otimes 1_k ) \circ C( C( V ) ) = ( \sigma_{d,d} \otimes \mathbb{I}_n ) \circ F( F( V ) ) = F( F( V ) )\circ ( \sigma_{d,d} \otimes \mathbb{I}_n ) = C( C( V ) ) \circ ( \sigma \otimes 1_k )$.
          Since $V$ was arbitrary, then $C( - )$ satisfies~\cref{Fig:ControlDefn:ContrlSym}.
    \item Let $V \in \CatUnitary_d( k, k )$ and $j \in \{ 1, 2, \ldots, d \}$.
          Since $\gamma_{k,j}$ is self-adjoint, then it follows by (\textbf{C4}) that $C( \gamma_{k,j} \circ V \circ \gamma_{k,j} ) = ( 1 \otimes \gamma_{k,j} ) \circ C( V ) \circ ( 1 \otimes \gamma_{k,j} )$.
          Since $V$ and $j$ were arbitrary, then $C( - )$ satisfies~\cref{Fig:ControlDefn:DataSym}.
    \end{itemize}
    This means that $C( - )$ is a control functor.
    Since every morphism in $\CatUnitary_d$ is unitary, then $C( - )$ is a dagger control functor by~\cref{Lem:UnitaryCtrl}.
    Finally, if $V, P \in \CatUnitary_d( k, k )$, then by (\textbf{C4}), $C( P^\dagger \circ V \circ P ) = F( P^\dagger \circ V \circ P ) = ( \mathbb{I}_d \otimes P^\dagger ) \circ F( V ) \circ ( \mathbb{I}_d \otimes P ) = ( 1 \otimes P^\dagger ) \circ C( V ) \circ ( 1 \otimes P )$.
    Since $V$ and $P$ were arbitrary, then $C( - )$ is a conjugated control functor.

    \vspace{1em}
    \noindent
    \textbf{Pointed Control Functors}.
    %
    Assume that $\ket{u} \in \mathbb{C}^d$ and $\ket{v} \in \mathbb{C}^d$.
    Since $U$ is Hermitian, then $\mathbb{C}^d$ can be partitioned into the $(+1)$ and $(-1)$ eigenspaces of $U$, where $\alpha_j = 0$ and $\alpha_j = 1$ respectively.
    To this end, define the following two sets where $[n] = \{ 1, 2, \ldots, n \}$.
    \begin{align*}
        X &= \{ b_j : j \in [d] \text{ and } \alpha_j = 0 \}
        &
        Y &= \{ b_j : j \in [d] \text{ and } \alpha_j = 1 \}
    \end{align*}
    It follows that $\{ b_j : j \in [d] \} = X \cup Y$.
    This partitioning will be used to determine the conditions under which $\ket{u}$ and $\ket{v}$ are points of $C( - )$.
    \begin{itemize}
    \item Assume that $C( - )$ is $( u, v )$-pointed.
          It will be shown that $U\ket{u} = -\ket{u}$.
          It then follows by a symmetric argument that $U\ket{v} = \ket{v}$.
          Construct a matrix $M = \sum_{j=0}^{d-1} \ket{d-j} \bra{j}$, so that $M$ is a permutation of the standard basis vectors within $\mathbb{C}^d$.
          Since $M$ is a permutation matrix, then $M$ is unitary, and hence $M \in \CatUnitary_d( 1, 1 )$.
          Since $C( - )$ is a $( u, v )$-pointed control functor, then the following equation holds where $\gamma_x = \braket{x|u}$.
          \begin{align*}
              \ket{u} \otimes M
              &=
              C( M ) \circ \left( \ket{u} \otimes \mathbb{I}_d \right)
              \\
              &=
              \sum_{j=1}^n \ket{b_j} \braket{b_j|u} \otimes M^{\alpha_j}
              &
              ( \text{by \cref{Thm:GenCtrlNF}} )
              \\
              &=
              \sum_{b_j \in X} \gamma_{b_j} \ket{b_j} \otimes M^0
              +
              \sum_{b_j \in Y} \gamma_{b_j} \ket{b_j} \otimes M^1
              \\
              &=
              \sum_{b_j \in X} \gamma_{b_j} \ket{b_j} \otimes \mathbb{I}_d
              +
              \sum_{b_j \in Y} \gamma_{b_j} \ket{b_j} \otimes M
          \end{align*}
          Since $\otimes$ is bilinear, then the following equation holds.
          \begin{equation*}
              \left( \ket{u} - \sum_{b_j \in Y} \gamma_{b_j} \ket{b_j} \right) \otimes M
              =
              \left( \sum_{b_j \in X} \gamma_{b_j} \ket{b_j} \right) \otimes \mathbb{I}_d
          \end{equation*}
          Since $\{ \ket{\ell} \bra{k} \}_{j=1,\ell=1}^{d}$ is a basis for $\CatUnitary_d( 1, 1 )$, then $\mathbb{I}_d$ and $M$ are not linear combinations of one-another.
          Then the left-hand side and right-hand side also not linear combinations of one-another.
          Then by equating components, $\ket{u} = \sum_{b_j \in Y} \gamma_{b_j} \ket{b_j}$.
          Then $\ket{u}$ is in the $(-1)$-eigenspace of $U$.
          In conclusion, $U\ket{u} = -\ket{u}$.
    \item Assume that $U\ket{u} = -\ket{u}$ and $U\ket{v} = \ket{v}$.
          Since $U\ket{u} = -\ket{u}$, then by definition $\ket{u}$ is in the $(-1)$-eigenspace of $U$.
          This means that $\ket{u} = \sum_{b_j \in Y} \gamma_j \ket{b_j}$ where $\gamma_j = \braket{b_j|u}$.
          Recall that $\braket{b_j|b_k} = \delta_{j,k}$ where $\delta$ is the Kronecker delta function.
          Let $V \in \CatUnitary_d( k ,k )$ and $n = d^k$.
          Then the following equation holds.
          \begin{align*}
              C( V ) \circ ( \ket{u} \otimes \mathbb{I}_n )
              &=
              \sum_{j=1}^n \ket{b_j} \braket{b_j|u} \otimes V^{\alpha_j}
              &
              ( \text{by \cref{Thm:GenCtrlNF}} )
              \\
              &=
              \sum_{j=1}^n \ket{b_j} \left( \sum_{b_k \in X} \gamma_k \braket{b_j|b_k} \right) \otimes V^{\alpha_j}
              &
              ( \text{by the bilinearity of } \braket{-|-} )
              \\
              &=
              \sum_{j=1}^n \ket{b_j} \left( \sum_{b_k \in X} \gamma_x \delta_{j,k} \right) \otimes V^{\alpha_j}
              \\
              &=
              \sum_{j=1}^n \, \sum_{b_k \in X} \gamma_k \delta_{j,k} \ket{b_j} \otimes V^{\alpha_j}
              &
              ( \text{by the bilinearity of } \otimes )
              \\
              &=
              \sum_{b_j \in X} \gamma_j \ket{b_j} \otimes V^1
              \\
              &=
              \sum_{b_j \in X} \gamma_j \ket{b_j} \otimes V
              \\
              &=
              \left( \sum_{b_j \in X} \gamma_j \ket{b_j} \right) \otimes V
              =
              \ket{u} \otimes V
              &
              ( \text{by the bilinearity of } \otimes )
          \end{align*}
          Since $V$ was arbitrary, then $C( V ) \circ ( \ket{u} \otimes \mathbb{I}_n ) = \ket{u} \otimes V$ for all $V \in \CatUnitary_d( k, k )$ with $n = d^k$.
          By a symmetric argument, starting from the fact that $\ket{v}$ in the $(+1)$-eigenspace of $U$, it follows that $C( V ) \circ ( \ket{v} \otimes \mathbb{I}_n ) = \ket{v} \otimes \mathbb{I}_n$ for all $V \in \CatUnitary_d( k, k )$ with $n = d^k$.
          Then $C( - )$ is $( u, v )$-pointed.
    \end{itemize}
    In conclusion, $C( - )$ is $( u, v )$-pointed if and only if $U\ket{u} = -\ket{u}$ and $U\ket{v} = \ket{v}$.
\end{proof}